\section{Introduction} 
% Current models good at coding tasks
% Agentic loops became more popular
% LLMs hallucinate by generating correct structure but false idea/logic
% Current hallucination detection include hidden states, token prob, predictive entropy, semantic uncertainty
% Current solutions often only detect but not fix problem
% No solutions with coding agentic loops
% This paper explores effect of Uncertainty Score on agents
Current Large Language Models (LLMs) have demonstrated outstanding performance
in various types of tasks, including question answering \cite{muludi2024retrieval}, summarization \cite{pu2023summarization}, and most
importantly software engineering \cite{zhang2026survey}. Integration of LLMs in agentic loops allowed 
to fully automate many processes and substantially increase efficiency and quality
of software engineers work.

Even so, truthfulness and quality of LLMs generations still suffers from hallucinations problem.
Hallucination refers to situation when model answers are structurally correct, but logically or 
factuality false. This may happen due to many hard to manage factors \cite{huang2025survey} such as poor quality of training
data, problems with context size or even corrupted context by user itself. Hallucinated
model output often leads to wrong conclusions, and misleading instructions,
that may confuse user. Agentic systems suffer from this even more - one poor design decision or incorrect tool
use may corrupt the whole context, heavily influencing future actions, and in most cases leading to poor software
quality or completely broken project.

Present papers manly focussed on hallucination detection \cite{azaria2023internalstatellmknows}, \cite{zhang2025icr}, \cite{shelmanov2025head}, \cite{fadeeva2024fact} in various ways of estimating 
model uncertainty. One of proposed approaches uses model output token probabilities to quantify 
uncertainty, hence possibly detection hallucinated sample. Another approaches use method of sampling
different responses on same input promt, focusing on estimating uncertainty of specific fact,
rather that some token. One of the most promising ways of uncertainty estimation is to use internal
model representations, collected from various layers, hence capturing hallucination, even before token generated.
Many approaches also explore combinations of proposed methods to train a separate simple MLP or Attention Head, that
classify false and true samples.

However, even with sufficient number of papers exploring hallucination detection, there are still limited
number of research actually fixing false generations \cite{andriushchenko2026efficient}, \cite{zhao-etal-2025-uncertainty}. Even so, in most cases this solutions only oriented
on simple Q/A tasks or one promt interaction with LLMs, but not complex long context interactions with models, that the case for coding agents.

This research aims to investigates the role of uncertainty score on agent behaviour by integrating various ways of
hallucination estimators directly into agentic loops. As a result we present a complete agentic system,
capable of detection and fixing hallucinated samples, hence increasing overall performance.

Main contribution of this paper is the following:
\begin{itemize}
    \item \textbf{Hallucination Aware Agentic Loop} - we propose hallucination aware agentic loop, that uses
    various uncertainty scores, for detecting and fixing hallucinated samples.
\end{itemize}
